\scanpage{49}
\mspage{12}
\[
(14)\qquad\alpha:\Delta_{r'}\longrightarrow\Delta_r
\]
est donnée par
\[
\alpha^*(x,i_*)=(x,\underbrace{\alpha^*(i_*)}_{\substack{\Vert\,\mathrm{def}\\i_*\circ\alpha}})
\]
i.e. elle est induite par la structure semi-
simpliciale produit dans le deuxième mb de
\[
(15)\qquad X_{\Delta_*}\subset X\times I_{\Delta_*}.
\]
\uline{Remarques}. a) Toutes ces constructions se généralisent au
cas où on remplacerait \inserted{l'ens. ordonné} $I$ par une catégorie \inserted{$\unsure{\mathcal{C}|C}$}, et
l'application $X\mapsto X_i$ par un foncteur $\unsure{\mathcal{C}|C}\longrightarrow\text{catégorie}$
des parties fermées de $X$ et de leurs inclusions
....

b) Si on part d'une famille $(X_i)_{i\in I}$, \inserted{$I$ un ens} sans
plus, on peut introduire ad hoc une relation
d'ordre dans $I$ par \deleted{$i\leq j\unsure{=|\Longleftrightarrow}$}
\[
i\leq j\ \overset{\mathrm{def}}{\Longleftrightarrow}\ X_i\leq X_j.
\]
Mais dans la présentation adoptée, il n'est pas
nécessaire de supposer que \inserted{l'appl. croissante} $I\longrightarrow\text{ens. des parties}$
fermées de $X$ soit injective, et encore moins
que ce soit un plongement d'ens. ordonnés,
i.e. que l'on ait
\[
X_i\leq X_j\ \Longrightarrow\ i\leq j
\]
qui est une réciproque de (1' b)).
\margin{48}
